GAMMA (Gobierno Automatizado del Maestro de Materiales) es una plataforma de asistencia inteligente que automatiza la normalización, detección de duplicados y clasificación de categoría en el flujo de creación de materiales sobre SAP, desarrollada para una unidad de negocio del sector energético con presencia en varios países de Centroamérica y el Caribe. El proyecto surge de un catálogo con 45,397 materiales distribuidos en un máximo de 1,538 categorías posibles, donde la ausencia de herramientas analíticas obligaba al gestor a validar manualmente cada solicitud, con un tiempo de gestión promedio de 3.3 horas por solicitud.

La plataforma integra tres módulos: un modelo de lenguaje para normalizar descripciones de texto libre, un motor de similitud difusa sobre PostgreSQL para detectar duplicados, y un clasificador supervisado —seleccionado tras evaluar comparativamente siete configuraciones— para sugerir la categoría correcta. El modelo ganador, una máquina de vectores de soporte lineal sobre representación de n-gramas de carácter, alcanzó una exactitud de 84.9,% en su primera sugerencia y 94.0,% considerando las tres principales.

La herramienta se validó en producción con el equipo de gestores entre el 18 de junio y el 21 de julio de 2026. La comparación contra la línea base mostró una reducción del tiempo de gestión de 58.4,% (p \textless\ 0.0001, prueba U de Mann-Whitney), y un retorno de inversión de 308x sobre el costo del proyecto. La encuesta de satisfacción aplicada al cierre del período confirmó una percepción favorable entre los gestores, aunque con una muestra reducida que limita su representatividad. Los resultados respaldan la viabilidad técnica y operativa de aplicar aprendizaje automático a la gobernanza de datos maestros en un entorno productivo real.